%% Supplementary Material for FAME manuscript
%% Target journal: Computers & Operations Research
\documentclass[a4paper,fleqn]{cas-sc}

\usepackage{amsmath,amssymb}
\usepackage{booktabs}
\usepackage{graphicx}
\usepackage{array}
\usepackage{multirow}
\usepackage{enumitem}
\usepackage{flafter}

% Supplement-specific numbering
\renewcommand{\thesection}{S\arabic{section}}
\renewcommand{\thesubsection}{S\arabic{section}.\arabic{subsection}}
\renewcommand{\thetable}{S\arabic{table}}
\renewcommand{\thefigure}{S\arabic{figure}}
\renewcommand{\theequation}{S\arabic{equation}}

\begin{document}

\shorttitle{Supplementary Material: FAME}
\shortauthors{Pereira and Oliveira}

% Keep tables/figures off the manual title area.
\suppressfloats[t]

\begin{center}
{\LARGE\bfseries Supplementary Material for\par}
\vspace{0.35em}
{\Large Temporal Transferability of Prediction-to-Decision Representations\\
in Constrained Resource Allocation\par}
\vspace{1.1em}
{\normalsize Tiago Martins Pereira$^{a,*}$ \quad Diana Campos de Oliveira$^{a}$\par}
\vspace{0.55em}
{\small\itshape
$^{a}$Statistics Department, Federal University of Ouro Preto (UFOP),\\
Campus Morro do Cruzeiro, Bauxita, Ouro Preto, 35402-136, Minas Gerais, Brazil\par}
\vspace{0.35em}
{\small $^{*}$Corresponding author: tiago.martin@ufop.edu.br\par}
\end{center}

\vspace{1.0em}
\hrule
\vspace{1.2em}

\section{Purpose and Scope}
\label{sec:supp_scope}

This Supplementary Material reports diagnostic and robustness analyses that
support the temporal transfer experiments presented in the main manuscript.
None of the analyses reported here changes the primary
development--calibration--freeze--evaluation protocol. In particular,
evaluation-period outcomes are never used to reselect the representation whose
future performance is reported in the main experiment.

The supplementary analyses address the robustness checks referenced in the
main manuscript. Section~S2 reports fantasy-football component ablation,
captain-selection sensitivity, target-convention and missing-next-market
sensitivity, calibration-surface stability, and moving-block bootstrap results.
Section~S3 reports the Energy operational robustness grid and temporal stability
of its calibrated representation. Section~S4 documents the reproducibility
audit trail, and Section~S5 states the interpretive boundaries of all
retrospective diagnostics.

\section{Fantasy-Football Robustness Analyses}
\label{sec:supp_fantasy}

\subsection{Component Ablation}
\label{sec:supp_ablation}

The primary frozen DOC representation for 2025 uses
$\widehat{\mathbf w}_{DOC}=(0.75,0.15,0.10)$ for expected performance,
upside/dispersion, and economic efficiency, respectively. To examine the
contribution of individual representation components, the 2025 evaluation was
repeated under prespecified ablations.

For the ``No upside'' and ``No economic efficiency'' variants, the corresponding
component is removed from the frozen DOC representation and the remaining DOC
weights are renormalized to sum to one. The ``No availability adjustment''
variant replaces the availability-adjusted expected component by the
conditional prediction and reconstructs the associated upside and economic
components consistently, while retaining the frozen DOC weights. The remaining
rows provide reference representations that isolate normalization, conditional
prediction, and equal weighting.

\begin{table}[!ht]
\centering
\caption{Fantasy-football component ablation in the out-of-time 2025
evaluation period. Utility is realized lineup utility without captain bonus.}
\label{tab:supp_ablation}
\begin{tabular}{lrr}
\toprule
Strategy & Mean utility & Cumulative utility \\
\midrule
Availability-adjusted expected prediction & 59.092 & 2186.40 \\
No upside & 56.853 & 2103.57 \\
Expected-only normalized representation & 56.759 & 2100.08 \\
Full FAME-DOC & 54.746 & 2025.61 \\
Equal-weight representation & 53.901 & 1994.34 \\
No economic efficiency & 53.860 & 1992.83 \\
Pure conditional prediction & 52.279 & 1934.32 \\
No availability adjustment & 48.417 & 1791.42 \\
\bottomrule
\end{tabular}
\end{table}

The ablation results show that the full three-component DOC representation is
not uniformly superior in the 2025 environment. Removing the upside component
raises mean utility from 54.746 to 56.853, whereas removing economic efficiency
reduces mean utility slightly to 53.860. Eliminating availability adjustment
produces the largest deterioration among the ablated FAME variants. These
results are diagnostic and do not constitute a second model-selection stage:
the representation evaluated prospectively remains the one frozen after 2024
DOC calibration.

\begin{figure}[!ht]
\centering
\includegraphics[width=.95\textwidth]{fig_2025_ablation.png}
\caption{Out-of-time 2025 component-ablation analysis.}
\label{fig:supp_ablation}
\end{figure}

\subsection{Captain-Selection Sensitivity}
\label{sec:supp_captain}

Captain assignment is evaluated as a secondary decision after the lineup has
already been selected. The experiment applies a 50\% captain bonus to the
realized score of the chosen non-coach player. Four captain rules are evaluated
for all lineup strategies: highest market value, highest historical mean,
highest expected performance, and highest conditional prediction. For FAME-DOC
and the original FAME representation, a fifth rule selects the player with the
highest corresponding FAME representation value.

\begin{table}[!ht]
\centering
\caption{Captain-rule sensitivity for the frozen FAME-DOC lineup in 2025.}
\label{tab:supp_captain}
\begin{tabular}{lrrr}
\toprule
Captain rule & Mean bonus & Mean final utility & Oracle hit rate \\
\midrule
Highest market value & 3.495 & 58.241 & 16.22\% \\
Highest historical mean & 3.330 & 58.076 & 16.22\% \\
Highest expected performance & 3.314 & 58.060 & 13.51\% \\
Highest conditional prediction & 2.861 & 57.607 & 5.41\% \\
Highest FAME representation & 2.730 & 57.476 & 16.22\% \\
\bottomrule
\end{tabular}
\end{table}

The choice of captain rule changes the level of final team utility but does not
alter the lineup generated by the frozen representation. Among the evaluated
FAME-DOC captain rules, highest market value produces the largest mean final
utility (58.241), while selecting the captain directly from the FAME
representation produces 57.476. The analysis therefore supports treating
captain selection as a distinct secondary decision rather than incorporating
evaluation-period captain performance into DOC recalibration.

\begin{figure}[!ht]
\centering
\includegraphics[width=.95\textwidth]{fig_2025_captain_selection.png}
\caption{Sensitivity of final team utility to alternative captain-selection
rules.}
\label{fig:supp_captain}
\end{figure}

\subsection{Target-Convention Audit and Missing-Next-Market Sensitivity}
\label{sec:supp_target}

The target-construction audit shows that a valid next-market observation is
available for more than 99.3\% of decision rows in every season from 2021 to
2025. The proportion of rows without a next-market observation ranges from
0.516\% to 0.623\%. Table~\ref{tab:supp_target_audit} summarizes the audit.

\begin{table}[!ht]
\centering
\caption{Target-alignment audit by season.}
\label{tab:supp_target_audit}
\begin{tabular}{lrrr}
\toprule
Season & Decision rows & Next market available & Missing next market \\
\midrule
2021 & 29,284 & 99.484\% & 0.516\% \\
2022 & 29,227 & 99.429\% & 0.571\% \\
2023 & 28,290 & 99.413\% & 0.587\% \\
2024 & 27,703 & 99.415\% & 0.585\% \\
2025 & 27,770 & 99.377\% & 0.623\% \\
\bottomrule
\end{tabular}
\end{table}

To assess sensitivity to the convention used for these uncommon missing
next-market rows, the 2025 operational analysis is repeated under two scenarios:
treating the missing next-market target as zero and excluding such observations.
The availability-adjusted expected-prediction reference is unchanged at 59.092
mean utility. FAME-DOC changes from 54.746 under the zero convention to 55.352
when the missing rows are excluded. The substantive conclusion of the primary
2025 comparison is therefore unchanged.

\begin{table}[!ht]
\centering
\caption{Sensitivity to the treatment of missing next-market observations.}
\label{tab:supp_missing_market}
\begin{tabular}{lrr}
\toprule
Missing-next-market convention & FAME-DOC & Expected prediction \\
\midrule
Missing target treated as zero & 54.746 & 59.092 \\
Rows with missing target excluded & 55.352 & 59.092 \\
\bottomrule
\end{tabular}
\end{table}

\subsection{Calibration-Surface Stability}
\label{sec:supp_stability}

DOC is optimized on a finite calibration sample, so the exact maximizing vector
need not be sharply identified. The 2024 DOC selection was therefore repeated
over 5,000 bootstrap resamples of calibration-period decision origins. The
original optimum, $(0.75,0.15,0.10)$, was selected in 1,200 resamples
(24.00\%). The expected-only vector $(1,0,0)$ was selected in 14.88\%, and
$(0.90,0.05,0.05)$ in 14.26\%. The distribution of bootstrap selections
therefore supports the interpretation of a favorable region rather than a
uniquely dominant representation.

\begin{table}[!ht]
\centering
\caption{Most frequently selected vectors in the 2024 DOC bootstrap
stability analysis.}
\label{tab:supp_bootstrap_top}
\begin{tabular}{lrr}
\toprule
Weight vector $(w_P,w_T,w_R)$ & Selection count & Frequency \\
\midrule
$(0.75,0.15,0.10)$ & 1200 & 24.00\% \\
$(1.00,0.00,0.00)$ & 744 & 14.88\% \\
$(0.90,0.05,0.05)$ & 713 & 14.26\% \\
$(0.80,0.15,0.05)$ & 448 & 8.96\% \\
$(0.85,0.00,0.15)$ & 333 & 6.66\% \\
\bottomrule
\end{tabular}
\end{table}

\begin{figure}[!ht]
\centering
\includegraphics[width=.95\textwidth]{fig_2024_doc_selection_stability.png}
\caption{Bootstrap selection frequencies for the most frequently selected
2024 DOC weight vectors.}
\label{fig:supp_doc_stability}
\end{figure}

A complementary local sensitivity analysis evaluates all weight vectors within
$L_1$ neighborhoods of the frozen 2024 solution. For radius 0.10, seven vectors
are included and their mean 2025 utility is 54.246, with a range from 52.813 to
55.814. For radius 0.20, 19 vectors are included and mean 2025 utility is
54.680, ranging from 51.995 to 56.801. Thus, the precise frozen maximizer is
only one member of a local region whose future performance is heterogeneous.

\begin{table}[!ht]
\centering
\caption{Local sensitivity around the frozen 2024 DOC vector.}
\label{tab:supp_local}
\begin{tabular}{rrrrr}
\toprule
$L_1$ radius & Vectors & Mean 2024 & Mean 2025 & 2025 range \\
\midrule
0.10 & 7  & 56.317 & 54.246 & 52.813--55.814 \\
0.20 & 19 & 56.051 & 54.680 & 51.995--56.801 \\
\bottomrule
\end{tabular}
\end{table}

\subsection{Calibration and Retrospective Utility Landscapes}
\label{sec:supp_landscapes}

The calibration and evaluation-period utility landscapes provide a geometric
diagnostic of representation transfer. Figure~\ref{fig:supp_calibration_landscape}
shows the 2024 DOC calibration surface from which the frozen vector was selected.
Figure~\ref{fig:supp_posthoc_landscape} shows the corresponding 2025 surface
computed retrospectively after the evaluation period. The latter is used only
to diagnose displacement of the favorable region and is never used to modify
the prospective 2025 decision rule.

\begin{figure}[!ht]
\centering
\includegraphics[width=.95\textwidth]{fig_2024_doc_calibration_landscape.png}
\caption{DOC calibration landscape in 2024. The selected representation is
frozen before the 2025 evaluation.}
\label{fig:supp_calibration_landscape}
\end{figure}

\begin{figure}[!ht]
\centering
\includegraphics[width=.95\textwidth]{fig_2025_posthoc_landscape.png}
\caption{Retrospective 2025 utility landscape. This diagnostic uses realized
2025 outcomes only after prospective evaluation and is not a deployable
selection rule.}
\label{fig:supp_posthoc_landscape}
\end{figure}

\subsection{Moving-Block Bootstrap}
\label{sec:supp_block_bootstrap}

The primary paired contrast contains 37 round-level observations and may exhibit
short-range temporal dependence. A circular moving-block bootstrap was therefore
computed for block lengths from two to six rounds. All resulting 95\% intervals
contain zero, matching the conclusion from the paired iid bootstrap.

\begin{table}[!ht]
\centering
\caption{Bootstrap sensitivity for FAME-DOC minus the normalized expected-only
representation in the primary 2024--2025 transition.}
\label{tab:supp_mbb}
\begin{tabular}{lrr}
\toprule
Method / block length & Lower 95\% limit & Upper 95\% limit \\
\midrule
Paired iid bootstrap & -6.253 & 1.924 \\
Moving block, 2 & -6.746 & 2.207 \\
Moving block, 3 & -6.963 & 2.410 \\
Moving block, 4 & -6.633 & 2.303 \\
Moving block, 5 & -6.715 & 2.300 \\
Moving block, 6 & -6.713 & 2.468 \\
\bottomrule
\end{tabular}
\end{table}

\section{Energy Robustness Analyses}
\label{sec:supp_energy}

\subsection{Operational Robustness Grid}
\label{sec:supp_energy_grid}

The Energy robustness analysis evaluates five capacity-margin settings and four
shortage-penalty multipliers,
\[
M\in\{1.00,1.10,1.20,1.30,1.40\},
\qquad
\lambda_{\mathrm{VOLL}}\in\{0.5,1,2,5\},
\]
for each of nine temporal replications and each predictive source. This yields
180 replication--condition combinations per predictive source. These operating
conditions are structured sensitivity settings, not statistically independent
observations.

\begin{table}[!ht]
\centering
\caption{Aggregate operational transfer across the Energy robustness grid.}
\label{tab:supp_energy_transfer}
\begin{tabular}{lrrrr}
\toprule
Source & Positive / total & Transfer rate & Mean gain & Minimum gain \\
\midrule
ENTSO-E & 152/180 & 84.44\% & 42.44\% & -60.30\% \\
SARIMAX & 180/180 & 100.00\% & 70.16\% & 3.26\% \\
LSTM    & 180/180 & 100.00\% & 66.47\% & 7.65\% \\
\bottomrule
\end{tabular}
\end{table}

SARIMAX and LSTM retain positive future DOC gains throughout the complete
robustness grid. ENTSO-E is less stable: 152 of 180 conditions remain positive,
and the minimum observed gain reaches -60.30\%. Because the robustness grid
changes the operational environment while retaining the chronological
calibration--freeze--evaluation logic, these results assess operational
sensitivity rather than provide additional independent replications.

\subsection{Temporal Stability of the Energy Representation}
\label{sec:supp_energy_theta}

The Energy representation is
\[
V_d^{(m)}(\theta)=\widehat D_d^{(m)}\exp(\theta).
\]
Across the nine temporal replications, mean calibrated $\widehat{\theta}$ is
0.085 for ENTSO-E, 0.132 for SARIMAX, and 0.108 for LSTM. The corresponding
temporal standard deviations are 0.054, 0.014, and 0.016, respectively.
ENTSO-E therefore exhibits substantially larger temporal variation in the
prediction-to-capacity transformation than the two internally estimated
forecast models.

\begin{table}[!ht]
\centering
\caption{Temporal stability of the calibrated Energy representation.}
\label{tab:supp_theta}
\begin{tabular}{lrr}
\toprule
Source & Mean $\widehat{\theta}$ & Temporal SD \\
\midrule
ENTSO-E & 0.085 & 0.054 \\
SARIMAX & 0.132 & 0.014 \\
LSTM & 0.108 & 0.016 \\
\bottomrule
\end{tabular}
\end{table}

The extended search domain
$\Theta=\{-0.30,-0.275,\ldots,0.275,0.30\}$
also provides a boundary diagnostic. The final calibrated solutions are
predominantly interior, so the reported transfer pattern is not explained by
systematic truncation at the upper calibration boundary.

\section{Reproducibility and Audit Trail}
\label{sec:supp_repro}

The experimental outputs retain a machine-readable audit trail spanning
predictive development, DOC calibration, frozen representations, future
decisions, and operational comparisons. In the fantasy-football pipeline, the
freeze audit confirms that predictive hyperparameters, ensemble weights, DOC
weights, budget and admissible formations were specified without using 2025
outcomes; 2025 is used only for sequential out-of-time evaluation.

The supplementary evidence package includes temporal-development folds,
predictive search candidates and selected hyperparameters, out-of-fold
predictions, walk-forward predictions, DOC calibration surfaces, bootstrap
selection frequencies, frozen DOC weights, round-level operational outputs,
post-hoc utility surfaces, ablation and captain analyses, target-convention
audits, and leakage/freeze audits. The Energy package analogously contains
temporal protocol definitions, calibration surfaces, frozen parameters,
replication-level results, and the complete operational robustness grid.

\section{Interpretive Boundaries}
\label{sec:supp_boundaries}

The analyses in this supplement are robustness and diagnostic analyses. They
should not be interpreted as additional opportunities to optimize the reported
2025 strategy after observing 2025 outcomes. In particular, the superior 2025
performance of some ablated or reference strategies is retrospective evidence
about representation transfer, not a replacement for the frozen DOC result.

Similarly, the retrospective utility landscape and local-neighborhood analyses
are used to characterize displacement of the favorable representation after
deployment. They are not implementable prospective policies. This distinction
preserves the information boundary required by the primary temporal
transferability claim.

\end{document}
